\documentclass[12pt]{article}
\usepackage{amsmath, amssymb}
\usepackage{geometry}
\usepackage{enumitem}
\geometry{margin=1in}

\title{CSE 400: Fundamentals of Probability in Computing\\
Tutorial--1 Solution}
\author{School of Engineering and Applied Science (SEAS)\\
Ahmedabad University}
\date{February 13, 2026}

\begin{document}

\maketitle

\section*{Definitions and Notation}

\subsection*{Multinomial Formula}
\[
\frac{n!}{n_1! n_2! \dots n_k!}
\]

\subsection*{Uniform Distribution}
Arrival time uniformly distributed over an interval.

\subsection*{Complement Rule}
\[
P(A^c) = 1 - P(A)
\]

\subsection*{Inclusion--Exclusion Principle}
\[
P(A \cup B) = P(A) + P(B) - P(A \cap B)
\]

\subsection*{De Morgan's Law}
\[
P(A^c \cap B^c) = 1 - P(A \cup B)
\]

\subsection*{Poisson Distribution}
If $X \sim \text{Poisson}(\lambda)$,
\[
P(X = k) = \frac{e^{-\lambda} \lambda^k}{k!}
\]

\subsection*{Gaussian (Normal) Distribution}
\[
f_X(x) = \frac{1}{\sqrt{2\pi\sigma^2}} 
\exp\left(-\frac{(x-m)^2}{2\sigma^2}\right)
\]

\[
F_X(x) = \Phi\left(\frac{x-m}{\sigma}\right),
\quad
P(X > x) = Q\left(\frac{x-m}{\sigma}\right)
\]

\subsection*{Binomial Distribution}
\[
P(X = k) = \binom{n}{k} p^k (1-p)^{n-k}
\]

\subsection*{Cumulative Distribution Function (Discrete)}
\[
F(a) = \sum_{x \le a} p(x)
\]

\section*{Assumptions}

\begin{itemize}
\item Order within groups does not matter and groups are unlabeled (Q1).
\item Arrival time is uniformly distributed over 30 minutes (Q2).
\item Independence of events where stated (Q8, Q10).
\item Poisson model applies for rare events in fixed intervals (Q4--Q6).
\item Jurors act independently with probability $\theta$ (Q8).
\item System components function independently with probability $p$ (Q10).
\end{itemize}

\section*{Theorems / Results}

\subsection*{Multinomial Counting Result}
\[
\frac{20!}{(5!)^4 4!}
\]

\subsection*{Poisson Complement}
\[
P(X \ge k) = 1 - \sum_{i=0}^{k-1} P(X=i)
\]

\subsection*{Gaussian Symmetry}
\[
\Phi(-x) = 1 - \Phi(x)
\]

\subsection*{Condition for Larger System Reliability}
\[
p > \frac{1}{2}
\]

\section*{Proofs / Derivations}

\subsection*{Q1}

\[
\frac{20!}{(5!)^4 4!}
\]

Probability:
\[
P = \frac{(5!)^4 4!}{20!}
\]

\subsection*{Q2}

Total interval = 30 minutes.

\[
P = \frac{10}{30} = \frac{1}{3}
\]

\subsection*{Q3}

Given:
\[
P(L)=0.15, \quad P(E)=0.25, \quad P(L \cap E)=0.08
\]

\[
P(E^c)=0.75
\]

\[
P(L \cup E)=0.15+0.25-0.08=0.32
\]

\[
P(L^c \cap E^c)=1-0.32=0.68
\]

\[
P(L^c | E^c)=\frac{0.68}{0.75}\approx 0.907
\]

\subsection*{Q4 ($\lambda=0.2$)}

\[
P(0)=e^{-0.2}=0.8187
\]

\[
P(1)=0.2 e^{-0.2}=0.1637
\]

\[
P(X \ge 2)=1-0.8187-0.1637=0.0176
\]

\subsection*{Q5 ($\lambda=3.5$)}

\[
P(0)=e^{-3.5}=0.0302
\]

\[
P(1)=3.5 e^{-3.5}=0.1057
\]

\[
P(X \ge 2)=0.8641
\]

\[
P(X \le 1)=0.1359
\]

\subsection*{Q6 ($\lambda=3$)}

\[
P(0)=0.0498
\]

\[
P(1)=0.1494
\]

\[
P(2)=0.2240
\]

\[
P(X \ge 3)=0.5768
\]

\[
P(X \ge 3 | X \ge 1)=\frac{0.5768}{0.9502}\approx 0.607
\]

\subsection*{Q7 ($m=-3, \sigma=2$)}

\[
P(X \le 0)=\Phi\left(\frac{3}{2}\right)
\]

\[
P(X > 4)=Q\left(\frac{7}{2}\right)
\]

\[
P(|X+3|<2)=2\Phi(1)-1
\]

\[
P(|X-2|>1)=1-Q(2)+Q(3)
\]

\subsection*{Q8 (Jury Problem)}

If innocent:
\[
\sum_{i=5}^{12} \binom{12}{i} \theta^i (1-\theta)^{12-i}
\]

If guilty:
\[
\sum_{i=8}^{12} \binom{12}{i} \theta^i (1-\theta)^{12-i}
\]

Overall:
\[
\alpha \sum_{i=8}^{12} \binom{12}{i} \theta^i (1-\theta)^{12-i}
+
(1-\alpha) \sum_{i=5}^{12} \binom{12}{i} \theta^i (1-\theta)^{12-i}
\]

\subsection*{Q9}

\[
c e^{\lambda}=1 \Rightarrow c=e^{-\lambda}
\]

\[
P(0)=e^{-\lambda}
\]

\[
P(X>2)=1-e^{-\lambda}-\lambda e^{-\lambda}-\frac{\lambda^2}{2}e^{-\lambda}
\]

\subsection*{Q10}

5-component system:
\[
10p^3(1-p)^2 + 5p^4(1-p) + p^5
\]

3-component system:
\[
3p^2(1-p)+p^3
\]

Condition:
\[
3(p-1)^2(2p-1) > 0
\]

\[
p > \frac{1}{2}
\]

\section*{Conclusion}

\begin{itemize}
\item Multinomial counting applied for equal partitioning.
\item Complement rule used for Poisson tail probabilities.
\item Conditional probability applied using complements and inclusion--exclusion.
\item Gaussian probabilities expressed using $\Phi$ and $Q$ functions.
\item Larger odd-component systems are better when $p > \frac{1}{2}$.
\end{itemize}

\end{document}